\chapter{The Solitonic Origin of Matter}
\label{chap:solitons}

Part~I took a single refusal---no singularities---and followed it to a cosmology.
The refusal was implemented by one object: a gravitational potential, extended by
torsion, that grows steeply enough at high density to halt collapse and reverse it
(Chapter~\ref{chap:bounce}). This Part takes that potential at its word in a second
place. If a torsion-supported potential is a real feature of geometry sourced by spin,
and not merely a cosmological convenience, then it does not switch on only when a
universe's worth of matter collapses---it is present wherever spin density is. And a
potential that can arrest the collapse of a universe can also, in principle, \emph{bind}:
it can hold a finite, self-sustaining knot of field together. The conjecture of this
Part is that those knots are the elementary fermions.

\begin{conjecture}
If the gravity--torsion potential is what forbids singularities, then its stable
bound solutions---regular, horizonless, spin-supported solitons---are candidates for
the elementary constituents of matter.
\end{conjecture}

This has a long pedigree (Wheeler's geons, the Soler and Friedberg--Lee soliton tradition,
and above all Heisenberg's nonlinear-spinor \emph{Weltformel}: one nonlinear Dirac equation,
one fundamental constant, the whole particle spectrum). It also has well-known walls. The
task here is to state the conjecture precisely, walk it to each wall, and report honestly
which fall, which stand, and which are merely untested. The status is held claim by claim;
the central quantitative prize---masses as binding energies rather than inputs---is named at
the end as the calculation the programme owes.

\section{The torsiton}
\label{sec:torsiton}

In Einstein--Cartan geometry the connection carries an antisymmetric part, torsion,
sourced by fermion spin (Axiom~\ref{ax:ec}). Integrating torsion out leaves a Dirac field
with a four-fermion self-interaction, the Hehl--Datta term,
\[
  \mathcal{L}_{\mathrm{int}} \;\sim\; -\frac{3}{2}\,\frac{8\pi G}{c^4}
  \left(\bar\psi\,\gamma_5\gamma^\mu\,\psi\right)^2 ,
\]
so the binding ``potential'' is not external: it is the well the fermion's own axial-spin
density digs for itself. Curvature pulls the field together; torsion, at high spin density,
pushes it apart. The same opposition that reverses a cosmological collapse admits a
\emph{stationary} balance---a localized lump held together permanently. We call such a
stationary, self-bound solution of the gravity--torsion field the \textbf{torsiton}: where
the gravity quantum is the graviton, the torsion-bound matter quantum is the torsiton, and
a torsiton is what an electron or a quark \emph{is}. (``Soliton'' remains the generic
mathematical object; ``torsiton'' names this particular one.)

That such a solution can exist in three dimensions, where Derrick's theorem forbids static
scalar lumps, is not an oversight: the first-order Dirac kinetic term scales differently
from a scalar's, and the binding balances two terms of different size-scaling---the same
two-terms-balance that makes the bounce. Soler-type nonlinear-Dirac solitons are known to
exist for exactly this reason. The self-consistent calculation confirms it: filling the
lowest levels (Pauli), sourcing the four-fermion field from their density, and iterating
the well to a fixed point yields a genuine self-bound object
(Figure~\ref{fig:self-consistent}). The gate is open---the torsiton exists, not as an
assumption but as the converged solution of the field's own equations.

\begin{figure}[t]
  \centering
  \includegraphics[width=\textwidth]{self_consistent}
  \caption{The self-consistent gravity--torsion soliton. The Hartree loop
  (levels~$\to$~density~$\to$~four-fermion field~$\to$~well) converges to a genuine
  self-bound profile: the torsiton exists as a solution, not a postulate. Reproducible
  from \texttt{sims/src/cartasis\_sims/self\_consistent.py}.}
  \label{fig:self-consistent}
\end{figure}

\section{Mass is configurational}
\label{sec:configurational-mass}

The decisive property follows at once. A torsiton's mass is not a label hung on a point;
it is the bound energy of a field configuration---roughly $94\%$ field energy, with the
rest the level energy of the trapped fermion. Mass is therefore \emph{configurational}: it
is set by how tightly the knot binds to the surrounding condensate, and it vanishes when the
condensate vanishes. The chiral-soliton calculation makes this explicit
(Figure~\ref{fig:chiral}): the soliton mass scales linearly with the condensate amplitude
$v$, so as $v \to 0$ the mass switches off.

\begin{figure}[t]
  \centering
  \includegraphics[width=\textwidth]{chiral_soliton}
  \caption{Configurational mass. The torsiton's mass is dominantly field energy and is
  proportional to the symmetry-breaking condensate $v=f_\pi$: turn off the condensate and
  the mass vanishes. Mass is a property of the bound configuration, not an intrinsic tag.}
  \label{fig:chiral}
\end{figure}

This is the bridge back to Part~I. What the bounce conserves are energy--momentum, charge,
spin, and $B\!-\!L$---never mass, which is configurational. So the bounce \emph{can} dissolve
matter: on the way into the membrane the condensate melts ($v\to0$), mass switches off, and
the field crosses as a hot, symmetric, effectively massless plasma; on the way out it
re-condenses and re-acquires weight. Catharsis (Chapter~\ref{chap:bounce}) and genesis
(Chapter~\ref{chap:genesis}) are the two halves of one crossing, and configurational mass is
why the crossing is allowed at all.

\section{Generations: an arithmetic ladder, a geometric hierarchy}
\label{sec:generations}

If the generations are internal excitation levels of the same torsiton---same charge,
colour, and isospin, different internal rung---their level energies are roughly evenly
spaced. But the observable mass is not the level energy; in the soliton picture it is the
overlap of the level's wavefunction with the localized, mass-giving substrate. Overlaps are
exponentially sensitive to the rung: a higher, more spread-out level overlaps the sharp
substrate exponentially less,
\[
  O_n \sim e^{-c\,n}, \qquad m_n \sim |O_n|^2 \sim e^{-2cn},
\]
so an \emph{arithmetic} ladder in $n$ becomes a \emph{geometric} ladder in mass. The
``unnaturally tiny'' electron Yukawa ($\sim 3\times 10^{-6}$) is then not a fine-tuned small
number but the exponential of an order-one overlap. The mechanism produces a large hierarchy
from $O(1)$ inputs---exactly the structure the Standard Model lacks---and a finite number of
rungs: the $S$-budget that binds the soliton runs out, capping the tower
(Figure~\ref{fig:gencap}), which is why there are three families and not thirty.

\begin{figure}[t]
  \centering
  \includegraphics[width=\textwidth]{generations_cap}
  \caption{Generations as a capped ladder. The internal excitation budget is finite, so the
  geometric mass tower terminates---a structural reason for a small, fixed number of
  generations rather than an arbitrary one.}
  \label{fig:gencap}
\end{figure}

How steep the geometric ladder is depends on the overlap with the \emph{substrate}---the
sharp chiral-restored core, not the soliton's own broad condensate. Overlapping each level
against that sharp substrate (its width fixed by the well's own healing length, not fitted)
turns a near-flat spread into a steep one of order $10^3$--$10^4$
(Figure~\ref{fig:substrate}), the observed size of the generation hierarchy. The exact ratios
and the per-tower steepness remain the chiral-dynamics refinement owed to the lattice; the
\emph{shape}---why generations span orders of magnitude---is the soliton's, with no fit.

\begin{figure}[t]
  \centering
  \includegraphics[width=\textwidth]{substrate_overlap}
  \caption{The substrate-overlap pass. Overlapping each generation with the sharply localized
  substrate (width $=$ the well's healing length, derived) rather than the soliton's broad
  self-condensate turns the spread from $\sim\!1$ into $\sim\!10^3$: the arithmetic level
  ladder becomes the observed geometric mass hierarchy, with no new fit.}
  \label{fig:substrate}
\end{figure}

\section{What is owed}

The single result that would make this Part worth more than its coherence is the recovery of
a fermion mass as a binding energy, parameter-free. That is not yet delivered: it is the
strongly-coupled, non-perturbative core of the theory, and it is named here as a calculation
to compute (Chapter~\ref{chap:weltformel} attempts the laptop-scale version; the
Appendix lists the lattice targets). What \emph{is} delivered is the structure---the torsiton
exists, mass is configurational, generations are a capped geometric ladder---and a clean
bridge to the cosmology: mass is exactly what the bounce does not conserve.
